\chapter{A Troubleshooting Method for the 28 Per Cent}\label{ch:tshootmethod}
\bp{2.4}

\begin{outcomes}
By the end of this chapter you will be able to:
\begin{tight}
  \item \textbf{Apply} a seven-step cycle to any fault, rather than guessing in a
        familiar order.
  \item \textbf{Choose} between the five structured methods and justify the
        choice from the symptom.
  \item \textbf{Scope} a fault before touching anything, using four questions
        that routinely halve the search.
  \item \textbf{Work} evidence-first: record before you change, change one thing,
        and undo what did not help.
\end{tight}
\end{outcomes}

\begin{prereq}
Everything. This chapter is the method that the rest of the book supplies the
knowledge for. \chref{transport} for the diagnostic tools and \chref{tshootl2}
for layer 2 in particular; the symptom index in Section~\ref{sec:symptomindex} points at all the
others.
\end{prereq}

\begin{keyterms}
scoping $\bullet$ top-down $\bullet$ bottom-up $\bullet$ divide and conquer
$\bullet$ follow the path $\bullet$ compare working to broken $\bullet$
hypothesis $\bullet$ blast radius $\bullet$ handover
\end{keyterms}

\begin{examnote}[What topic 2.4 actually says]
\emph{Troubleshoot basic Layer 2/Layer 3 connectivity and device operations using
show commands (including show logs), ping, extended ping, trace route, and packet
capture output.}

Every tool in that list has already been taught --- \chref{transport} for ping,
extended ping, traceroute and reading captures; \chref{management} for
\cmd{show logging}; \chref{tshootl2} for layer 2. This chapter deliberately does
not repeat them.

What it adds is the thing the tools cannot supply: \textbf{an order to use them
in}. Eight of the twenty-nine topics begin \emph{Troubleshoot} or
\emph{Diagnose}, which is 28\% of the blueprint, and in every one of them the
examiner already knows you can type \cmd{show ip interface brief}. What is being
tested is whether you can decide to.
\end{examnote}

\section{Why method beats knowledge under pressure}

Two engineers with identical knowledge produce very different results on a fault
neither has seen. The difference is almost never what they know; it is whether
they have an order to work in when the obvious answer turns out to be wrong.

An unstructured approach checks the thing that was broken last time, then the
thing the person before you suggested, then something at random. It works
quickly when the fault is familiar and collapses when it is not --- and the exam,
like a bad morning, specialises in faults that are not.

\begin{alertbox}[The failure mode to design against]
The characteristic mistake is \textbf{finding a fault and stopping}. There is a
misconfiguration on the path, you fix it, and the symptom persists --- because it
was not \emph{the} fault, or not the only one.

A method protects against this by making you state what you expect the fix to do
\emph{before} you apply it. If the prediction fails, you have learned something;
if you had no prediction, you have only changed something.
\end{alertbox}

\section{The cycle}

\ccnafigh{%
\begin{tikzpicture}[font=\scriptsize]
  \tikzset{
    st/.style={rounded corners=3pt, draw=#1, fill=#1!8, text=#1!45!black,
               line width=0.9pt, text width=2.55cm, align=center,
               minimum height=1.35cm, font=\tiny\bfseries},
    a/.style={-{Stealth[length=2.2mm]}, line width=0.9pt, neutral!70}}

  \node[st=dom1] (s1) at (0,0)     {1. DEFINE\\[2pt]\mdseries what exactly is
      broken, for whom};
  \node[st=dom1] (s2) at (3.35,0)  {2. GATHER\\[2pt]\mdseries scope it before
      touching anything};
  \node[st=dom1] (s3) at (6.7,0)   {3. ANALYSE\\[2pt]\mdseries which layers can
      this be};
  \node[st=dom1] (s4) at (10.05,0) {4. ELIMINATE\\[2pt]\mdseries rule out whole
      areas at once};
  \node[st=dom4] (s5) at (13.4,0)  {5. HYPOTHESISE\\[2pt]\mdseries one cause,
      and a prediction};

  \node[st=dom4] (s6) at (10.05,-2.3) {6. TEST\\[2pt]\mdseries change one thing,
      check the prediction};
  \node[st=dom2] (s7) at (13.4,-2.3)  {7. SOLVE\\[2pt]\mdseries fix, verify,
      \textbf{document}};

  \draw[a] (s1) -- (s2);  \draw[a] (s2) -- (s3);
  \draw[a] (s3) -- (s4);  \draw[a] (s4) -- (s5);
  \draw[a] (s5) -- (s6);
  \draw[a] (s6) -- (s7);

  % The label rides on its own arc, on a white ground. Beside the arc it either
  % lands in the gap the arc encloses -- which the ELIMINATE box occupies, so it
  % printed straight through "rule out whole areas at once" and neither could be
  % read -- or sits close enough to the grey 5-to-6 arrow to look like its label.
  \draw[a, dom4!70] (s6) to[bend left=30]
      node[font=\tiny, text=dom4!45!black, align=center, pos=0.34,
           fill=white, inner sep=1.5pt] {prediction\\failed} (s5);
\end{tikzpicture}}%
{The loop is between 5 and 6, and it is the point. A hypothesis that does not
predict correctly is discarded, not patched.}{tscycle}

\begin{conceptbox}[Step 5 is where amateurs and professionals separate]
``I think it is the trunk allowed list'' is not a hypothesis. \textbf{``I think
VLAN 30 is missing from the allowed list on the uplink, in which case
\cmd{show interfaces trunk} will not list it, and VLANs 10 and 20 will be
unaffected''} is one --- because it can be wrong, and you will know immediately.

Say what you expect to see before you run the command. It takes three seconds and
it converts every check into information whichever way it goes.
\end{conceptbox}

\section{Five structured methods}

\begin{longtable}{@{}L{3.2cm}L{5.4cm}L{6.4cm}@{}}
\toprule
\rowcolor{primary!15}
\textbf{Method} & \textbf{How it works} & \textbf{Use it when} \\
\midrule
\endhead
\textbf{Bottom-up} &
Start at layer 1 and work up. Cable, link light, duplex, VLAN, address, route &
You suspect hardware or physical change; a new install; nothing works at all \\
\rowcolor{lightbg}
\textbf{Top-down} &
Start at the application and work down &
Most of the network is fine and one application is not. Avoids inspecting a
switch for a DNS fault \\
\textbf{Divide and conquer} &
Start in the middle --- usually layer 3. Ping the gateway; if it works, layers
1--3 are fine &
You have no strong prior. \textbf{This is the default}, and the fastest on
average \\
\rowcolor{lightbg}
\textbf{Follow the path} &
Trace the actual route the traffic takes, device by device, checking each hop &
Connectivity fails between two specific points and works elsewhere \\
\textbf{Compare working to broken} &
Diff a failing device or port against an identical working one &
Some members of an identical set work. Frequently the fastest of all \\
\bottomrule
\end{longtable}

\begin{alertbox}[Compare-to-working is under-used and enormously effective]
Twenty-four ports configured the same way, one of them broken. Do not reason
about it --- \cmd{show running-config interface} on both and read the difference.

The same applies to a device: if three hypervisor hosts work and the fourth does
not (\chref{virtualisation}), the fault is in what is \emph{not} identical, and
that is a much smaller search space than ``what could be wrong with a hypervisor
uplink''.

Whenever a working comparison exists, use it first. It converts an open-ended
diagnosis into a diff.
\end{alertbox}

\section{Scope before you touch}

Four questions, asked before any command, routinely eliminate most of the
possible causes.

\begin{longtable}{@{}L{4.5cm}L{10.5cm}@{}}
\toprule
\rowcolor{primary!15}
\textbf{Question} & \textbf{What the answer eliminates} \\
\midrule
\endhead
\textbf{One user, or many?} &
One user: the host, its port, its cable. Many on one switch: that switch or its
uplink. Many across sites: routing, DNS, DHCP, or the service itself. This single
question sets the whole search space \\
\rowcolor{lightbg}
\textbf{When did it start, and has it ever worked?} &
``Never worked'' is a configuration fault --- something was built wrong.
``Worked until Tuesday'' is a change fault. They are searched completely
differently \\
\textbf{What changed?} &
Almost every fault in a previously working network follows a change. Ask about
changes to anything, not just the network --- a server move, a firmware update,
a contractor in the comms room \\
\rowcolor{lightbg}
\textbf{Does it fail by name but work by address?} &
Answers ``is this DNS'' in one test, and DNS is a large fraction of reported
network faults that are not network faults (\chref{dns}) \\
\bottomrule
\end{longtable}

\begin{conceptbox}[``What changed?'' is the highest-yield question in operations]
A network that worked on Monday and fails on Tuesday was changed. The change may
be undocumented, may have been made by someone else, may not look related, and
the person who made it may sincerely believe it could not have caused this.

Ask anyway, ask broadly, and check the evidence rather than the recollection:
\cmd{show logging} for \cmd{\%SYS-5-CONFIG\_I} messages (\chref{management}),
\cmd{show archive config differences}, and the change record.

The answer is often ``nothing'', and it is often wrong.
\end{conceptbox}

\section{Evidence discipline}

Three rules. They are boring, they are what separates a twenty-minute fault from
a three-hour one, and they are equally true in an exam simulation.

\begin{tight}
  \item \textbf{Record before you change.} Capture the current state ---
        \cmd{show running-config}, the relevant \cmd{show} output, the exact
        symptom. A configuration you overwrote without reading is evidence you
        destroyed, and if the change makes things worse you cannot get back.
  \item \textbf{Change one thing at a time.} Two simultaneous changes that fix
        the fault leave you not knowing which one did, and one of them may have
        introduced tomorrow's problem.
  \item \textbf{Undo what did not help.} An accumulation of speculative changes
        is how a network becomes undiagnosable. If your prediction was wrong,
        put it back before trying the next idea.
\end{tight}

\begin{alertbox}[The exception, and know that you are taking it]
Under a major outage there is pressure to restore service first and understand
later. That is a legitimate decision --- but take it deliberately, say out loud
that you are taking it, and capture what you can as you go. The failure is not
choosing speed; it is choosing speed by accident and having no idea afterwards
what you did.
\end{alertbox}

\section{Symptom index}\label{sec:symptomindex}

Where to go in this volume, given what the user actually said.

\begin{longtable}{@{}L{6.5cm}L{8.5cm}@{}}
\toprule
\rowcolor{primary!15}
\textbf{Symptom} & \textbf{Start here} \\
\midrule
\endhead
No link light, or errors climbing &
\chref{physical} --- cable, duplex, speed, transceiver \\
\rowcolor{lightbg}
Host has \cmd{169.254.x.x} &
\chref{dhcp}, then \chref{tshootl2}. Nothing answered its request \\
Works within its VLAN, nothing beyond &
Gateway or mask: \chref{clienttshoot}, then \chref{routingtable} \\
\rowcolor{lightbg}
Works by address, fails by name &
\chref{dns}. Not a network fault \\
One whole VLAN dead, others fine &
\chref{vlans}, \chref{trunking} --- allowed list, SVI, or a per-SVI setting \\
\rowcolor{lightbg}
Intermittent, slow, MAC flapping &
\chref{stp} --- suspect a loop \\
Port \cmd{err-disabled} &
\chref{l2security} for port security or BPDU guard; \chref{management} to read
the reason \\
\rowcolor{lightbg}
Some destinations reachable, others not &
\chref{routingtable} for longest-prefix match, then \chref{acl} \\
OSPF neighbours will not form &
\chref{ospfv2} --- area, timers, MTU, network type, duplicate router ID \\
\rowcolor{lightbg}
Internet works for a while then stops &
\chref{nat} --- pool exhaustion, or a missing \cmd{overload} \\
Two gateways both claim to be active &
\chref{fhrp} --- group numbers or subnet mismatch \\
\rowcolor{lightbg}
Wireless: connects, no access &
\chref{clienttshoot} --- security parameters, or MAC randomisation \\
Wireless: strong signal, unusable &
\chref{wireless} --- check the noise floor, not the signal \\
\rowcolor{lightbg}
VMs on one host unreachable &
\chref{virtualisation} --- the uplink is probably an access port \\
``The log shows nothing'' &
\chref{management} --- check the severity thresholds first \\
\bottomrule
\end{longtable}

\section{Handing it over}

Faults outlive shifts. A handover that says ``still looking at the VLAN issue''
wastes the next person's first hour.

\begin{tight}
  \item \textbf{The symptom}, precisely, and who is affected.
  \item \textbf{What has been ruled out}, and the evidence that ruled it out.
        This is the most valuable part and the most often omitted.
  \item \textbf{What has been changed}, including anything not yet undone.
  \item \textbf{The current hypothesis} and the next test.
\end{tight}

\begin{pitfall}
\begin{tight}
  \item \textbf{Starting to type before scoping.} Four questions first.
  \item \textbf{Checking what broke last time.} Familiarity is not evidence.
  \item \textbf{No prediction before a change.} You learn nothing from the
        result.
  \item \textbf{Two changes at once.} Now you do not know which worked.
  \item \textbf{Leaving failed changes in place.} They become the next fault.
  \item \textbf{Fixing the first fault found and declaring victory} without
        re-testing the original symptom.
  \item \textbf{Ignoring an available working comparison.}
  \item \textbf{Believing ``nothing changed''.} Check the logs.
  \item \textbf{Not writing it down.} You will meet this fault again.
\end{tight}
\end{pitfall}

\begin{breakfix}[Two faults at once, and fixing the first one changes nothing.]
Users on the second floor report that the file server is unreachable. You find
that the uplink from \cmd{SW2} has VLAN~20 missing from its trunk allowed list,
which would certainly do it. You add it, confirm \cmd{show interfaces trunk} now
lists VLAN~20 --- and the users report no improvement at all.

\textbf{The instinct} is to assume the fix did not take, and to re-check the
trunk. Resist it: you verified it, and it is correct. A correct fix that does not
resolve the symptom means \textbf{your hypothesis explained a real fault, but not
this one}.

\textbf{Go back to scoping}, which is where the answer was all along.

\begin{tight}
  \item \textbf{One user or many?} All of the second floor --- consistent with
        the trunk.
  \item \textbf{Has anyone \emph{else} lost the server?} Nobody asked. It turns
        out the third floor cannot reach it either, and they are on a different
        switch with a correct trunk.
\end{tight}

That second answer eliminates the trunk as the cause of the reported symptom
immediately: a fault common to two switches with different configurations is not
in either switch.

\begin{verify}{R1}{show ip route 10.20.0.0}
% Network not in table
\end{verify}

\textbf{Diagnosis.} The server's subnet is not in the routing table. Its route was
lost overnight --- and \emph{that} is why nobody can reach it, from any floor.

The missing VLAN on \cmd{SW2}'s trunk was a second, genuine, unrelated fault. It
had been there for weeks without being noticed, because nothing in VLAN~20 on the
second floor had needed the uplink until the server was investigated. Finding it
was useful. Believing it was the cause was the error.

\textbf{The lesson, and it is the one this chapter exists for.} Two things went
wrong here, and only one of them was reported:

\begin{tight}
  \item \textbf{The scope was never established.} One question --- ``is anyone
        else affected?'' --- would have ruled out the trunk before it was
        touched.
  \item \textbf{There was no prediction.} ``Adding VLAN 20 will restore access
        for the second floor'' is testable, and it failed. Because no prediction
        was stated, the failure read as ``the fix did not work'' rather than as
        ``the hypothesis was wrong'', and the next twenty minutes went into
        re-checking a trunk that was already correct.
\end{tight}

Fix the routing fault. Keep the trunk fix --- it was a real defect --- but record
it separately, as a second issue found during the investigation, so the change
record reflects what actually happened.
\end{breakfix}

\begin{lab}{Faults you did not create}{Packet Tracer or CML, in pairs}
\textbf{Platform note.} The value here is entirely in \emph{not} knowing the
fault, so this lab must be done in pairs or by exchanging saved files. Working
alone on a fault you injected yourself teaches nothing.

\textbf{Topology.} Any multi-VLAN, multi-router topology from Parts II and III
that both of you have working.

\begin{tasks}
  \item Each of you takes a working copy and injects \textbf{three} faults from
        different layers --- one physical or layer 2, one layer 3, one a service
        (DHCP, DNS, NAT or an ACL). Record them privately. At least one should
        be silent, producing no log message.
  \item Exchange files. Before running a single command, write down your scoping
        questions and what each possible answer would eliminate.
  \item Diagnose. For every check, write the prediction \emph{before} you run it,
        and whether it held. Keep this record --- it is the deliverable.
  \item Fix one fault at a time and re-test the original symptom after each. Note
        any case where a correct fix did not change the symptom, and what you
        concluded.
  \item Time yourself. Compare against your partner on the same set of faults.
  \item Write the handover you would give if you ran out of time halfway.
  \item \textbf{Round 2.} Inject three more, but this time include \textbf{two
        faults on the same path} so that fixing one changes nothing. This is the
        Break \& Fix above, and it is the most valuable version of this exercise.
  \item \textbf{Round 3.} Inject one fault into \emph{one} of twenty-four
        identically configured ports. Time how long it takes with reasoning
        versus with a configuration diff.
  \item \textbf{Extension.} Build a one-page fault-finding card for your own use:
        the four scoping questions, the five methods, and the symptom index
        trimmed to what you personally forget.
\end{tasks}
\end{lab}

\begin{checkpoint}
\begin{tightnum}
  \item Give the seven steps of the cycle and say which two form a loop.
  \item Name the five structured methods and give the symptom that selects each.
  \item What are the four scoping questions, and what does each eliminate?
  \item Why must you state a prediction before making a change?
  \item A correct fix does not resolve the reported symptom. What does this tell
        you, and what should you do next?
  \item What are the four things a handover must contain?
\end{tightnum}
\end{checkpoint}

\begin{chaptersummary}
\begin{tight}
  \item Define, gather, analyse, eliminate, hypothesise, test, solve --- and loop
        between hypothesis and test until the prediction holds.
  \item A hypothesis must predict something that could be wrong. State it before
        you run the command.
  \item Five methods: bottom-up for physical, top-down for one application,
        divide and conquer as the default, follow the path between two points,
        and compare-to-working whenever an identical working case exists.
  \item Scope first: one user or many, has it ever worked, what changed, does it
        fail by name but work by address.
  \item ``What changed'' is the highest-yield question, and ``nothing'' is often
        wrong. Check \cmd{show logging} and the archive diff.
  \item Record before changing, change one thing, undo what did not help.
  \item A correct fix that does not resolve the symptom means the hypothesis was
        wrong --- or that there is more than one fault.
  \item A handover states the symptom, what has been ruled out and how, what was
        changed, and the next test.
\end{tight}
\end{chaptersummary}

\begin{reviewq}
  \item \textbf{(MCQ)} Which method should you reach for first when eight of
        twenty-four identically configured ports are failing?
        \begin{tight}
          \item[(a)] Bottom-up \quad (b)~Top-down
          \item[(c)] Compare working to broken \quad (d)~Follow the path
        \end{tight}
  \item \textbf{(MCQ)} A user reports the intranet is down. \cmd{ping} by IP
        address succeeds and by name fails. Which method has this already
        applied, and what does the result establish?
        \begin{tight}
          \item[(a)] Bottom-up; the cable is fine
          \item[(b)] Divide and conquer; layers 1--3 work and the fault is DNS
          \item[(c)] Follow the path; the route is correct
          \item[(d)] Compare to working; the host is misconfigured
        \end{tight}
  \item \textbf{(Short)} Explain why ``I think it is the trunk'' is not a
        hypothesis, and rewrite it as one.
  \item \textbf{(Short)} You fix a genuine misconfiguration and the reported
        symptom is unchanged. Give the two possible explanations and the next
        step for each.
  \item \textbf{(Applied)} A whole floor loses connectivity at 09:10 on a Monday.
        Write out your first five minutes: the questions you ask, the commands
        you run, and what each result would eliminate. Do not assume a cause.
  \item \textbf{(Applied)} You must hand over an unresolved fault at the end of a
        shift. Write the handover, inventing plausible detail, and mark which
        part of it the next engineer will find most valuable and why.
\end{reviewq}
